\section{Fakta dan Aturan}

\textbf{Fakta} adalah predikat yang menyatakan kebenaran yang selalu berlaku tanpa syarat \cite{learn_prolog,ucsd_prolog,spivey_prolog}. Sintaksnya berupa nama predikat diikuti argumen dalam tanda kurung, ditutup titik: \texttt{predikat(arg1, arg2, ...).} Fakta tidak mengandung variabel yang perlu diisi; setiap argumen berupa konstanta (atom atau bilangan). Contoh \texttt{orangtua(ahmad, ali).} menyatakan bahwa Ahmad adalah orang tua dari Ali. Arity predikat---jumlah argumennya---menentukan struktur relasi; predikat \texttt{orangtua/2} bersifat biner karena menghubungkan dua entitas \cite{ucsd_prolog,rose_prolog}.

\textbf{Aturan} memiliki bentuk \texttt{Kepala :- Tubuh.} dan dibaca ``Kepala benar jika Tubuh benar'' \cite{learn_prolog,swi_manual,lu_mead_prolog}. Simbol \texttt{:-} menyatakan implikasi; Kepala adalah satu literal (atom predikat), sedangkan Tubuh dapat berisi konjungsi literal yang dipisahkan koma. Setiap koma bermakna ``dan'' (AND): semua kondisi di Tubuh harus terpenuhi agar Kepala dapat disimpulkan. Misalnya \texttt{ayah(X,Y) :- orangtua(X,Y), laki(X).} menyatakan X adalah ayah dari Y jika X adalah orang tua Y dan X berjenis kelamin laki-laki \cite{learn_prolog,lu_mead_prolog}.

Tabel~\ref{tab:sintaks-prolog} merangkum bentuk sintaks fakta, aturan, dan query beserta maknanya \cite{learn_prolog,swi_manual}.

\begin{table}[htbp]
\centering
\small
\begin{tabular}{|l|>{\raggedright\arraybackslash}p{3.85cm}|>{\raggedright\arraybackslash}p{4.15cm}|>{\raggedright\arraybackslash}p{4.15cm}|}
\hline
\textbf{Komponen} & \textbf{Bentuk} & \textbf{Contoh} & \textbf{Makna} \\
\hline
Fakta & \texttt{predikat(arg1,...)}. & \texttt{orangtua(ahmad, ali).} & Predikat selalu benar \\
\hline
Aturan & \texttt{Kepala :- Tubuh.} & \texttt{ayah(X,Y) :- orangtua(X,Y), laki(X).} & Kepala benar jika Tubuh benar \\
\hline
Query & \texttt{?- goal.} & \texttt{?- kakek(X, ali).} & Tujuan yang ingin dibuktikan \\
\hline
\end{tabular}
\caption{Ringkasan sintaks fakta, aturan, dan query Prolog \cite{learn_prolog,swi_manual}}
\label{tab:sintaks-prolog}
\end{table}

Dalam logika, fakta dan aturan Prolog bersesuaian dengan \logic{klausa Horn} \cite{learn_prolog,olp_fol,endriss_prolog}. Klausa Horn adalah implikasi dengan paling banyak satu literal positif; dalam Prolog, kepala aturan adalah literal positif tunggal, dan tubuh adalah konjungsi literal (yang dapat bernilai negatif dalam ekstensi tertentu). Fakta adalah klausa Horn tanpa tubuh---implikasi dengan konklusi saja. Keterbatasan pada klausa Horn memungkinkan Prolog menggunakan \emph{SLD resolution} (Selective Linear Definite clause resolution) sebagai mesin inferensi yang efisien \cite{learn_prolog,olp_fol}.

Dua pandangan semantik berlaku bagi program Prolog \cite{learn_prolog,spivey_prolog}. \emph{Semantik deklaratif} memandang program sebagai teori logika: program benar secara logis jika dan hanya jika fakta dan aturan mengimplikasikan query. \emph{Semantik prosedural} memandang cara Prolog mengevaluasi: urutan klausa dan urutan literal dalam tubuh mempengaruhi urutan pencarian dan efisiensi. Pemula disarankan memikirkan program secara deklaratif terlebih dahulu; pemahaman prosedural berguna untuk optimasi dan debugging \cite{learn_prolog,ucsd_prolog}.

Predikat dapat bersifat unary (satu argumen, misalnya \texttt{laki(ahmad)}) atau binary serta higher-arity \cite{ucsd_prolog,rose_prolog}. Konvensi penamaan: gunakan huruf kecil untuk atom dan predikat, huruf besar untuk variabel. Predikat silsilah keluarga sering memakai arity 2: \texttt{orangtua(X,Y)}, \texttt{ayah(X,Y)}, \texttt{kakek(X,Z)}. Aturan
\begin{quote}\texttt{kakek(X,Z) :- ayah(X,Y), orangtua(Y,Z).}\end{quote}
mendefinisikan kakek sebagai ayah dari orang tua; variabel Y menghubungkan dua generasi \cite{learn_prolog,swi_manual}.

\begin{contoh}
\begin{lstlisting}[style=prologstyle]
orangtua(ahmad, ali).
orangtua(ahmad, sari).
laki(ahmad).
ayah(X, Y) :- orangtua(X, Y), laki(X).
kakek(X, Z) :- ayah(X, Y), orangtua(Y, Z).
\end{lstlisting}
\end{contoh}

\begin{contoh}
Fakta dan aturan untuk relasi ``teman'' dan ``saudara''. Dua orang saudara jika memiliki orang tua yang sama; dua orang teman jika keduanya saudara atau didefinisikan eksplisit \cite{learn_prolog}.
\begin{lstlisting}[style=prologstyle, caption={Predikat teman dan saudara}]
orangtua(ahmad, ali).
orangtua(ahmad, sari).
orangtua(ahmad, budi).
saudara(X, Y) :- orangtua(P, X), orangtua(P, Y), X %*\=*) Y.
teman(ali, budi).
teman(sari, budi).
teman(X, Y) :- saudara(X, Y).
\end{lstlisting}
Query \texttt{?- saudara(ali, sari).} menjawab \texttt{true}. Query \texttt{?- teman(X, budi).} dengan menekan \texttt{;} dapat mengembalikan \texttt{X = ali}, lalu \texttt{X = sari}, lalu \texttt{X = ali} lagi (melalui aturan saudara).
\end{contoh}

\begin{konsep}
Fakta: predikat yang selalu benar, tanpa variabel bebas. Aturan: \texttt{Kepala :- Tubuh} berarti Kepala jika Tubuh; koma di Tubuh = konjungsi (AND). Fakta dan aturan bersesuaian dengan klausa Horn. Semantik deklaratif: apa yang benar; semantik prosedural: bagaimana Prolog membuktikan.
\end{konsep}
